\subsection{准素分解的唯一性性质}

定义 3. 设 M 是诺特环上的模，N 是 M 的子模。M 中 N 的准素分解 \(N = \bigcap_{i \in I} Q_i\) 称为既约的，如果满足下列条件：

\begin{enumerate}
\def\labelenumi{(\alph{enumi})}
\item
  不存在指标 \(i \in \mathbf{I}\)，使得 \(\bigcap_{j \neq i} Q_j \subset Q_i\)；
\item
  若 \(\operatorname{Ass}(\mathbf{M} / \mathbf{Q}_i) = \{\mathfrak{p}_i\}\)，则 \(\mathfrak{p}_i\)（\(i \in \mathbf{I}\)）彼此不同。
\end{enumerate}

从\(\mathbf{N} = \bigcap_{i\in \mathbf{I}}\mathbf{Q}_i\)在\(\mathbf{N}\)中的\(\mathbf{M}\)中可以如下导出\(\mathbf{M}\)在\(\mathbf{N}\)中：令\(\mathbf{J}\)为\(\mathbf{I}'\)的子集所成集合中满足\(\mathbf{I}\)使得\(\mathbf{N} = \bigcap_{i\in \mathbf{I}'}\mathbf{Q}_i\)。显然，\((\mathbf{Q}_i)_{i\in \mathbf{J}}\)满足条件 (a)。再令\(\Phi\)为由\(p_i\)组成且满足\(i\in J\)；对于所有\(p\in \Phi\)，令\(H(p)\)为由\(i\in J\)组成且满足\(p_i = p\)，并令\(Q(p) = \bigcap_{i\in H(p)}Q_i\)；由第 1 节命题 2 可知，\(Q(p)\)是\(p\)-准素的（相对于\(M\)）；此外\(N = \bigcap_{p\in \Phi}Q(p)\)，所以族\(Q((p))_{p\in \Phi}\)从而是\(N\)在\(M\)中的既约准素分解。我们将看到，第 2 节定理 1 的证明中定义的准素分解是既约的；这由下述命题推出：

命题 4. 设 M 是诺特环上的模，N 是 M 的子模，\(N = \bigcap_{i \in I} Q_i\) 是 N 在 M 中的准素分解；对于所有 \(i \in I\)，令 \(\{p_i\} = \text{Ass}(M/Q_i)\)。该分解为既约的充要条件是：各 \(p_i\) 彼此不同且属于 \(\text{Ass}(M/N)\)；此时

\begin{enumerate}
\def\labelenumi{(\arabic{enumi})}
\setcounter{enumi}{1}
\tightlist
\item
\end{enumerate}

\[
\operatorname{Ass} (\mathbf {M} / \mathbf {N}) = \bigcup_ {i \in I} \left\{\mathfrak {p} _ {i} \right\}\tag{3}
\]

\[
\operatorname{Ass} \left(\mathrm{Q} _ {i} / \mathrm{N}\right) = \bigcup_ {j \neq i} \left\{\mathfrak {p} _ {j} \right\} \quad \text{对所有 } i \in \mathrm{I}.
\]

若陈述中的条件满足，\(N = \bigcap_{j \neq i} Q_j\)不可能成立，因为否则可推出\(\text{Ass}(M/N) \subset \bigcup_{j \neq i} \{p_j\}\)（§ 1，第 1 节，命题 3 的推论 2），这与假设矛盾；于是 N 的准素分解\((Q_i)_{i \in I}\)必然是既约的。反之，\(\text{Ass}(M/N) \subset \bigcup_{i \in I} \{p_i\}\)总是成立（§ 1，第 1 节，命题 3 的推论 2）；另一方面，对所有\(i \in I\)，令\(P_i = \bigcap_{j \neq i} Q_j\)；则\(P_i \cap Q_i = N\)，且若\(P_i \neq N\)是既约的，则\((Q_i)_{i \in I}\)非零，并同构于子模\(P_i/N\)中的\((P_i + Q_i)/Q_i\)，从而\(M/Q_i\)（§ 1，第 1 节，命题 3 及命题 2 的推论 1）；由于\(\{p_i\} = \text{Ass}(P_i/N)\)，\(P_i/N \subset M/N\)，\(p_i \in \text{Ass}(M/N)\)这就完成了陈述中条件之必要性的证明以及公式 (2) 的证明。最后，由于\(N = \bigcap_{j \neq i} (Q_j \cap Q_i)\)，\(\text{Ass}(Q_i/N) \subset \bigcup_{j \neq i} \text{Ass}(Q_i/(Q_j \cap Q_i))\)（§ 1，第 1 节，命题 3 的推论 2）；但\(Q_i/(Q_j \cap Q_i)\)同构于子模\((Q_i + Q_j)/Q_j\)中的\(M/Q_j\)，故\(\text{Ass}(Q_i/(Q_j \cap Q_i)) \subset \{p_j\}\)，且并且

\[
\operatorname{Ass} \left(Q _ {i} / N\right) \subset \bigcup_ {j \neq i} \left\{\mathfrak {p} _ {j} \right\};
\]

计及 (2) 及 §1 第 1 节的命题 3，公式 (3) 即易得。

推论。设 A 是诺特环，M 是 A-模，N 是 M 的子模，且 \((\mathbf{Q}_{i})_{i\in I}\) 是 N 在 M 中的准素分解。则 \(\operatorname{Card}(\mathbf{I}) \geqslant \operatorname{Card}(\operatorname{Ass}(\mathbf{M}/\mathbf{N}))\)；要使 \((\mathbf{Q}_{i})_{i\in I}\) 成为既约准素分解，充要条件是 \(\operatorname{Card}(\mathbf{I}) = \operatorname{Card}(\operatorname{Ass}(\mathbf{M}/\mathbf{N}))\)。

由命题 4 前的说明可知，存在 \((\mathbf{R}_j)_{j\in J}\) 作为 \(N\) 在 \(M\) 中的既约准素分解，使得 \(\operatorname{Card}(J) \leqslant \operatorname{Card}(I)\)；第一个断言于是由第二个断言得出，而后者是命题 4 的结果。

命题 5. 设 A 是诺特环，M 是 A-模，N 是 M 的子模，\(\mathbf{N} = \bigcap_{i \in I} \mathbf{Q}_i\) 是 N 在 M 中的既约准素分解；对于所有 \(i \in I\)，令 \(\{\mathfrak{p}_i\} = \operatorname{Ass}(\mathbf{M} / \mathbf{Q}_i)\)。若 \(\mathfrak{p}_i\) 是 \(\operatorname{Ass}(\mathbf{M} / \mathbf{N})\) 的极小元，则 \(\mathbf{Q}_i\) 等于 N 关于 \(\mathfrak{p}_i\) 的饱和（第二章，第 2 节，第 4 节）（参见习题 2）。

显然可以只考虑 N = 0 的情形，必要时以 M/N 替换 M。若 \(p_{i}\) 是 Ass(M) 中的极小元，则 Ass(M) 中不与 A - \(p_{i}\) 相交的元素所成的集合归结为 \(p_{i}\)；命题于是由上面的公式 (3) 以及 §1，第 2 节，命题 6 得出，因为典范映射 \(M \to M_{p_{i}}\) 的核等于关于 \(p_{i}\) 的 0 的饱和（第二章，第 2 节，第 4 节）。

注。属于\(\mathfrak{p}_i\in\mathrm{Ass}(\mathbf{M}/\mathbf{N})\)的素理想，如果不是该集合的极小元，有时称为与\(\mathbf{M}/\mathbf{N}\)相伴的浸入素理想；若\(\mathbf{M}/\mathbf{N}\)是有限生成的，则对于\(\mathfrak{p}_0\in\mathrm{Ass}(\mathbf{M}/\mathbf{N})\)成为浸入素理想的充要条件是\(\mathrm{V}(\mathfrak{p}_0)\)不是\(\mathrm{Supp}(\mathbf{M}/\mathbf{N})\)的不可约分支（第二章，第 4 节，第 3 节，命题 14 的推论 2）；若\((\mathrm{Q}(\mathfrak{p}))_{\mathfrak{p}\in\mathrm{Ass}(\mathbf{M}/\mathbf{N})}\)和\((\mathrm{Q}'(\mathfrak{p}))_{\mathfrak{p}\in\mathrm{Ass}(\mathbf{M}/\mathbf{N})}\)是\(\mathbf{N}\)在\(\mathbf{M}\)中的两个既约准素分解，则可能有\(\mathrm{Q}'(\mathfrak{p}_0)\neq\mathrm{Q}(\mathfrak{p}_0)\)（习题 24 (c)）；总可以通过对其中出现的准素子模附加补充条件，定义\(\mathbf{N}\)在\(\mathbf{M}\)中的典范既约准素分解（习题 4）。

